% §5 Exercises
%
% Cross-reference resolution status:
%
% DEFINED here:
%   \label{ex:ch5-1} through \label{ex:ch5-11}
%
% RESOLVES forward references:
%   \Cref{ex:ch5-5} -- referenced from bloom-from-scratch.tex line 261
%                      (the optimal-k uniqueness exercise that §5.2's
%                      proof of \Cref{thm:bloom-optimal-k} defers to).
%   \Cref{ex:ch5-7} -- referenced from bloom-from-scratch.tex line 208
%                      (the bit-check independence counter-example that
%                      §5.2's FPR derivation defers to).
%
% RESOLVED (used here, defined elsewhere):
%   \Cref{thm:bloom-optimal-k}      -- §5.2 (Plan 5 Task 3)
%   \Cref{eq:bloom-fpr}             -- §5.2 (Plan 5 Task 3)
%   \Cref{cor:kronecker-set}        -- §5.3 (Plan 5 Task 4)
%   \Cref{sec:predicate-algebra}    -- §5.3 (Plan 5 Task 4)
%   \Cref{thm:kronecker}            -- chapter 4 §4.4 (Plan 1)
%   \Cref{thm:composition}          -- chapter 3 §3.3 (Plan 3)
%
% Difficulty-marker macros (not in alex.sty; defined locally if not yet defined):
\providecommand{\exstar}{$\star$}
\providecommand{\exstarstar}{$\star\star$}
\providecommand{\exstarstarstar}{$\star\star\star$}

\section*{Exercises}
\addcontentsline{toc}{section}{Exercises}

\begin{enumerate}

% -----------------------------------------------------------------------
% WARM-UP  (★ × 3): exercises 1, 2, 3
% -----------------------------------------------------------------------

\item \exstar
\label{ex:ch5-1}
\textbf{Predicate-to-rates.}
Let $\Universe = \{1, 2, \ldots, 10\}$ and $A = \{2, 5, 7\}$. An
approximate set $\ASet{A}$ is queried once on every element of
$\Universe$ and produces the following table of reports:
\[
  \begin{array}{c|cccccccccc}
    x & 1 & 2 & 3 & 4 & 5 & 6 & 7 & 8 & 9 & 10 \\\hline
    \ASet{A}(x) & \bot & \top & \top & \bot & \top & \bot & \bot & \bot & \top & \bot
  \end{array}
\]
\begin{enumerate}
  \item[(a)] Identify which queries hit in-set elements and which hit
    out-of-set elements; tabulate the four outcome counts (TP, FP, FN,
    TN) using $A$ as ground truth.
  \item[(b)] Compute the empirical false-positive rate $\hat{\fprate}$
    as the fraction of out-of-set queries that report $\top$, and the
    empirical false-negative rate $\hat{\fnrate}$ as the fraction of
    in-set queries that report $\bot$. Be explicit about the
    denominators.
  \item[(c)] State in one sentence why these two empirical rates are
    the natural sample-based estimators of the per-block parameters
    $(\fprate, \fnrate)$ that \Cref{sec:predicate-algebra} attaches to
    $\ASet{A}$.
\end{enumerate}

\item \exstar
\label{ex:ch5-2}
\textbf{Bloom filter parameters.}
A Bloom filter is built with $m = 8000$ bits, $n = 1000$ inserted
elements, and $k = 5$ hash functions.
\begin{enumerate}
  \item[(a)] Compute $\fprate$ from \Cref{eq:bloom-fpr}, that is,
    $\fprate \approx (1 - e^{-kn/m})^k$ with the given $(m, n, k)$.
    Report at least three significant figures.
  \item[(b)] For the same $(m, n)$, compute the optimal hash count
    $k^* = (m/n)\ln 2$ from \Cref{thm:bloom-optimal-k}. Round to the
    nearest integer if necessary and report both the real-valued
    optimum and the integer used.
  \item[(c)] Evaluate \eqref{eq:bloom-fpr} at $k = k^*$ and compare to
    the rate from (a). By what factor does the optimal $k$ reduce the
    false-positive rate at this $(m, n)$?
\end{enumerate}

\item \exstar
\label{ex:ch5-3}
\textbf{One-sided check.}
For each of the following data structures, state whether it is
one-sided in $\fnrate$ ($\fnrate = 0$, $\fprate > 0$), one-sided in
$\fprate$ ($\fprate = 0$, $\fnrate > 0$), two-sided (both rates
positive), or exact (both rates zero). Give a one-sentence
justification per item drawn from the construction.
\begin{enumerate}
  \item[(a)] A Bloom filter (\Cref{def:bloom-filter}).
  \item[(b)] An exact set-membership test on a hash table that stores
    every element of $A$ verbatim.
  \item[(c)] A cardinality estimator such as HyperLogLog used as a
    membership oracle by thresholding the estimated cardinality.
\end{enumerate}

% -----------------------------------------------------------------------
% MEDIUM  (★★ × 5): exercises 4, 5, 6, 7, 8
% -----------------------------------------------------------------------

\item \exstarstar
\label{ex:ch5-4}
\textbf{Notebook-tied Bloom filter empirics.}
Open the chapter 5 companion notebook
\texttt{python/notebooks/ch05-approximate-sets.ipynb}. Build a Bloom
filter with $m = 10000$ bits and $n = 1000$ inserted random strings,
and vary the hash count over $k \in \{1, 2, 3, \ldots, 15\}$.
\begin{enumerate}
  \item[(a)] For each $k$, estimate the empirical false-positive rate
    $\hat{\fprate}(k)$ by querying a large pool of strings known to be
    absent from the inserted set and recording the fraction reported
    as $\top$.
  \item[(b)] On the same axes, plot $\hat{\fprate}(k)$ and the
    theoretical curve $(1 - e^{-kn/m})^k$ from \Cref{eq:bloom-fpr}.
    Describe in one sentence how closely the two curves agree across
    the range of $k$.
  \item[(c)] Read off the integer $k$ that minimizes $\hat{\fprate}$
    and compare to the optimum $k^* = (m/n)\ln 2$ predicted by
    \Cref{thm:bloom-optimal-k}. Do the empirical and theoretical
    optima agree to within the one-bin resolution of your sweep?
\end{enumerate}

\item \exstarstar
\label{ex:ch5-5}
\textbf{Uniqueness of the optimal $k$.}
The proof of \Cref{thm:bloom-optimal-k} in \Cref{sec:bloom-from-scratch}
reduced the first-order optimality condition to
\[
  (1 - u)\ln(1 - u) \;-\; u \ln u \;=\; 0, \qquad u \in (0, 1),
\]
under the change of variable $u = e^{-kn/m}$, and observed by
inspection that $u = 1/2$ solves it. The proof then claimed that
$u = 1/2$ is the \emph{only} solution in $(0, 1)$, leaving the
uniqueness to this exercise.
\begin{enumerate}
  \item[(a)] Define $\phi(u) = (1 - u)\ln(1 - u) - u \ln u$ on
    $(0, 1)$. Verify the symmetry $\phi(1 - u) = -\phi(u)$, and
    confirm that $\phi(u) \to 0$ as $u \to 0^+$ and as $u \to 1^-$.
    Conclude that $\phi(1/2) = 0$ without further calculation.
  \item[(b)] Compute $\phi'(u)$ and verify that $\phi'(u) =
    -\ln(1 - u) - \ln u - 2$. Show that $\phi'$ attains its minimum
    at $u = 1/2$, where $\phi'(1/2) = 2 \ln 2 - 2 < 0$, and tends to
    $+\infty$ at both endpoints. Conclude that $\phi'$ changes sign
    exactly twice on $(0, 1)$, once on $(0, 1/2)$ and once on
    $(1/2, 1)$.
  \item[(c)] Use the sign pattern of (b) and the boundary behavior of
    (a) to conclude that $\phi$ is positive throughout $(0, 1/2)$
    and negative throughout $(1/2, 1)$, with the unique zero located
    at $u = 1/2$.
  \item[(d)] Transfer the conclusion back to $k$ via the change of
    variable $u = e^{-kn/m}$. Confirm that $k^* = (m/n) \ln 2$ is
    the unique interior critical point of $f(k) = (1 - e^{-kn/m})^k$
    on $(0, \infty)$, and combine with the second-derivative check
    in \Cref{sec:bloom-from-scratch} to conclude that it is the
    global minimum.
\end{enumerate}

\item \exstarstar
\label{ex:ch5-6}
\textbf{Predicate algebra: joint probability via Kronecker.}
Let $\Universe$ have $|\Universe| = 6$ elements and $A \subseteq
\Universe$ with $|A| = 2$. An approximate set $\ASet{A}$ has per-block
rates $\fprate = 0.1$ and $\fnrate = 0.05$ in the deterministic-rates
regime of \Cref{sec:predicate-algebra}.
\begin{enumerate}
  \item[(a)] Fix three specific out-of-set elements $x_1, x_2, x_3 \in
    \Universe \setminus A$. Compute the probability that $\ASet{A}$
    reports $\top$ on each of $x_1, x_2, x_3$ (three simultaneous
    false positives), assuming Axioms 1 and 2 are in force.
  \item[(b)] Write the same probability as a single entry of the
    Kronecker power $Q^{\otimes 3}$, identifying which row and column
    of the per-query matrix
    $Q = \bigl(\begin{smallmatrix} 0.95 & 0.05 \\ 0.10 & 0.90
    \end{smallmatrix}\bigr)$ each of the three factors reads.
  \item[(c)] Justify the factorization in one or two sentences by
    naming the role of \Cref{thm:kronecker} (which gives the joint
    matrix as a Kronecker product) and \Cref{cor:kronecker-set}
    (which specializes the product to the two-parameter approximate-set
    setting).
\end{enumerate}

\item \exstarstar
\label{ex:ch5-7}
\textbf{Bit-check independence counter-example.}
The derivation of $\fprate \approx (1 - e^{-kn/m})^k$ in
\Cref{sec:bloom-from-scratch} treats the $k$ bit-check events as
independent. \Cref{eq:bloom-fpr} is therefore an approximation; the
exact false-positive rate is in general different. This exercise
constructs the smallest worked counter-example.
\begin{enumerate}
  \item[(a)] Take $m = 4$ bits, $n = 1$ inserted element, $k = 2$
    independent hash functions, each uniform over $\{1, 2, 3, 4\}$.
    Enumerate all $4 \times 4 = 16$ assignments of $(h_1(\text{inserted}),
    h_2(\text{inserted}))$ with equal probability $1/16$; for each
    assignment, write out which bits are set.
  \item[(b)] Now consider a query on a fresh out-of-set element $y$.
    Its two hashes $(h_1(y), h_2(y))$ are also uniform over
    $\{1, \ldots, 4\}^2$, independent of $h_i(\text{inserted})$.
    Compute the exact probability that both $b[h_1(y)] = 1$ and
    $b[h_2(y)] = 1$, that is, the exact false-positive rate, by
    averaging over all $16 \times 16 = 256$ joint hash assignments.
  \item[(c)] Compare the exact rate from (b) to the approximation
    \eqref{eq:bloom-fpr}, $\fprate \approx (1 - e^{-2 \cdot 1/4})^2
    = (1 - e^{-1/2})^2 \approx 0.155$. State in one or two sentences
    why the two values differ and identify the source of the
    correlation: if $h_1(y) = h_2(y)$, the two ``independent''
    bit-check events collapse into one.
\end{enumerate}

\item \exstarstar
\label{ex:ch5-8}
\textbf{Two-sided vs one-sided space trade-off via the entropy floor.}
\Cref{sec:one-sided-two-sided} previews chapter~8's entropy floor: any
FNR-zero approximate set at false-positive rate $\fprate$ requires at
least $-\log_2 \fprate$ bits per stored element. The standard Bloom
filter, at its optimal $k^*$ (\Cref{thm:bloom-optimal-k}), pays
$\approx 1.44 \cdot \log_2(1/\fprate)$ bits per element.
\begin{enumerate}
  \item[(a)] Compute the entropy floor $-\log_2 \fprate$ in bits per
    element for the three target rates $\fprate \in \{0.1, 0.01,
    0.001\}$. Tabulate to one decimal place.
  \item[(b)] Compute the Bloom filter's actual bit-per-element cost
    $1.44 \cdot \log_2(1/\fprate)$ for the same three targets.
  \item[(c)] Report the absolute overhead in bits per element (Bloom
    cost minus entropy floor) at each target. Interpret in one
    sentence: as $\fprate$ shrinks, does the absolute overhead grow,
    shrink, or stay constant, and what does that mean for the
    practical cost of one-sidedness at very low false-positive rates?
\end{enumerate}

% -----------------------------------------------------------------------
% CHALLENGING / EXTENSION  (★★★ × 3): exercises 9, 10, 11
% -----------------------------------------------------------------------

\item \exstarstarstar
\label{ex:ch5-9}
\textbf{Counting Bloom filter sketch.}
The standard Bloom filter (\Cref{def:bloom-filter}) supports no
deletions: clearing a bit on \textsc{Delete}$(x)$ would clear bits
shared by other inserted elements and break their queries. A
\emph{counting Bloom filter} resolves this by replacing each bit with
a small unsigned integer counter (typically 4 bits); \textsc{Insert}
increments the $k$ chosen counters and \textsc{Delete} decrements
them; \textsc{Query} returns $\top$ if all $k$ counters are positive.
\begin{enumerate}
  \item[(a)] Argue that \textsc{Delete} is correct \emph{provided}
    no counter ever overflows during \textsc{Insert} and no counter
    is ever decremented for an element that was not previously
    inserted. Both conditions can fail in practice; identify the
    failure mode of each.
  \item[(b)] Sketch why a counting Bloom filter is in general
    two-sided rather than one-sided: counter overflow (saturation
    at the maximum value) and the use of small counters together
    introduce a positive false-negative rate. Counter overflow caps
    the count, so subsequent decrements may zero a counter that
    should still be positive, causing a query on an inserted element
    to return $\bot$.
  \item[(c)] Compare the space cost. A $b$-bit counter Bloom filter
    occupies $bm$ bits, compared with $m$ bits for the standard
    filter. For the typical choice $b = 4$, this is a $4\times$
    overhead. Comment in one or two sentences on the trade-off: the
    counting filter buys an additional API (\textsc{Delete}) at the
    price of $4\times$ space, structural two-sidedness, and a small
    overflow-induced loss.
\end{enumerate}

\item \exstarstarstar
\label{ex:ch5-10}
\textbf{Block-stability under non-uniform priors.}
\Cref{prop:block-stability} of \Cref{sec:predicate-algebra} invokes
Axiom 2 to collapse the per-element rates $\fprate(x), \fnrate(x)$
to per-block rates $\fprate, \fnrate$. The collapse depends on the
assumption that the random rate attached to each block is shared by
every element in that block. Construct a scenario where this
assumption fails.
\begin{enumerate}
  \item[(a)] Take a learned classifier (a noisy approximate-membership
    oracle, in our framework) whose accuracy varies across the
    out-of-set block: queries on out-of-set elements with very low
    prior probability of confusion are answered correctly with rate
    $1 - \fprate_{\text{easy}} = 0.99$, while queries on
    near-the-boundary out-of-set elements achieve only
    $1 - \fprate_{\text{hard}} = 0.70$. Describe how this construction
    gives a $\fprate(x)$ that genuinely depends on $x$ within the
    out-of-set block $B_2 = \Universe \setminus A$.
  \item[(b)] Argue informally that the construction violates the
    spirit of Axiom 2 even though the classifier itself is
    deterministic. (\emph{Hint:} Axiom 2 says the per-block rate is
    a single random variable shared by every block member; here the
    per-element rates are distinct, so no single block rate fits.)
  \item[(c)] State the immediate consequence: the Kronecker
    factorization of \Cref{cor:kronecker-set} no longer applies in
    its two-parameter form; the joint distribution over $m$ queries
    must instead carry per-element rates and is no longer
    universe-size-independent. Briefly identify which workloads make
    the gap visible and which mask it (\emph{hint:} a workload that
    samples uniformly from $B_2$ averages over the heterogeneity and
    looks like a single-rate workload; a workload that adversarially
    selects near-boundary out-of-set elements exposes the gap).
\end{enumerate}

\item \exstarstarstar
\label{ex:ch5-11}
\textbf{Forward to chapter 6: union FPR via the composition theorem.}
Chapter~6 will derive the false-positive rate of the union of two
approximate sets $\ASet{A}_1 \cup \ASet{A}_2$ from the rates of the
constituents. For two FNR-zero approximate sets with FPRs
$\fprate_1, \fprate_2$, both representing the same underlying set $A$
on the same universe, and assuming independent error events, the
union (implemented by bit-OR on Bloom filters with disjoint hash
families) has FPR
\[
  \fprate_{\text{union}} \;=\; 1 - (1 - \fprate_1)(1 - \fprate_2).
\]
\begin{enumerate}
  \item[(a)] Derive this expression from \Cref{thm:composition} of
    chapter~3. (\emph{Hint:} for a fixed out-of-set element $x$, the
    union reports $\top$ if and only if at least one of the two
    constituent queries reports $\top$. The complement is that both
    report $\bot$, which under independence factors as $(1 -
    \fprate_1)(1 - \fprate_2)$.)
  \item[(b)] Verify the formula on a numerical example: take
    $\fprate_1 = 0.10$ and $\fprate_2 = 0.05$ and compute
    $\fprate_{\text{union}}$. Confirm that the result lies between
    $\max(\fprate_1, \fprate_2)$ and $\fprate_1 + \fprate_2$, with
    the second-order term $\fprate_1 \fprate_2$ accounting for the
    gap.
  \item[(c)] Describe a concrete situation where the independence
    assumption breaks and the formula no longer applies. For
    instance, the two filters share the same hash family over the
    same universe; then the bit-OR event on the same out-of-set
    element $x$ is perfectly positively correlated across the two
    filters whenever both filters set the same bits, and the union
    FPR collapses to a single-filter FPR rather than the
    composition-theorem prediction. State the qualitative effect on
    the union rate (over-estimate or under-estimate) and identify
    which structural property of the construction the formula
    presupposes.
\end{enumerate}

\end{enumerate}
